\chapter{Gastric Banding Surgery}

\section*{What Is This Test or Procedure?}
Gastric banding, commonly called Lap-Band surgery, places an adjustable band around the upper stomach to create a small pouch and limit food intake.

\section*{Why This Test or Procedure Is Ordered}
It may be considered as a bariatric treatment for selected patients with severe obesity when structured lifestyle and medical treatment have not produced adequate benefit.

\section*{How This Test or Procedure Is Performed}
Usually laparoscopically, an adjustable band is placed around the upper stomach and connected by tubing to a port beneath the skin. The band can be adjusted by adding or removing saline.

\section*{Risks and Follow-up}
Slippage, pouch enlargement, reflux, vomiting, swallowing difficulty, infection, erosion, inadequate weight loss, and revision surgery are possible. Long-term nutritional, behavioral, and surgical follow-up is required.

\section*{References}
\begin{enumerate}
  \item MedlinePlus. Bariatric Surgery. U.S. National Library of Medicine; 2025.
  \item Current Medical Diagnosis and Treatment. McGraw-Hill; 2025.
\end{enumerate}
\clearpage
